% D7
% Auto-generated from markdown source
% origin: ai

\documentclass[12pt, a4paper]{article}
\usepackage[utf8]{inputenc}
\usepackage[T1]{fontenc}
\usepackage{amsmath, amssymb, amsfonts, amsthm}
\usepackage{textcomp}
\usepackage{geometry}
\usepackage{hyperref}
\usepackage{booktabs}
\usepackage{graphicx}
\usepackage{fancyhdr}
\usepackage{tcolorbox}
\usepackage{enumitem}

\geometry{margin=1in}
\pagestyle{fancy}
\fancyhf{}
\rhead{D7}
\lhead{\thepage}

\newtheorem{definition}{Definition}[section]
\newtheorem{theorem}{Theorem}[section]
\newtheorem{corollary}{Corollary}[section]
\newtheorem{proposition}{Proposition}[section]
\newtheorem{lemma}{Lemma}[section]
\newtheorem{conjecture}{Conjecture}[section]

\newcommand{\calI}{\mathcal{I}}
\newcommand{\calP}{\mathcal{P}}
\newcommand{\calE}{\mathcal{E}}
\newcommand{\calS}{\mathcal{S}}
\newcommand{\calH}{\mathcal{H}}
\newcommand{\calO}{\mathcal{O}}
\newcommand{\calF}{\mathcal{F}}
\newcommand{\calU}{\mathcal{U}}
\newcommand{\calV}{\mathcal{V}}
\newcommand{\calA}{\mathcal{A}}
\newcommand{\calG}{\mathcal{G}}
\newcommand{\calB}{\mathcal{B}}
\newcommand{\calC}{\mathcal{C}}
\newcommand{\calL}{\mathcal{L}}
\newcommand{\calM}{\mathcal{M}}
\newcommand{\calX}{\mathcal{X}}
\newcommand{\calY}{\mathcal{Y}}
\newcommand{\calW}{\mathcal{W}}
\newcommand{\calT}{\mathcal{T}}
\newcommand{\calZ}{\mathcal{Z}}
\newcommand{\calR}{\mathcal{R}}
\newcommand{\calQ}{\mathcal{Q}}
\newcommand{\calD}{\mathcal{D}}
\newcommand{\calN}{\mathcal{N}}
\newcommand{\calK}{\mathcal{K}}

\newcommand{\Tr}{\text{Tr}}
\newcommand{\E}{\mathbb{E}}
\newcommand{\diff}{\mathrm{d}}
\newcommand{\KL}{D_{\mathrm{KL}}}
\newcommand{\argmax}{\arg\max}
\newcommand{\argmin}{\arg\min}
\newcommand{\hbarF}{\hbar}

\newtcolorbox{thmbox}[1]{
  colback=blue!5!white, colframe=blue!75!black, fonttitle=\bfseries,
  title={#1}, breakable
}
\newtcolorbox{defbox}[1]{
  colback=green!5!white, colframe=green!75!black, fonttitle=\bfseries,
  title={#1}, breakable
}

\begin{document}

\title{D7}
\author{RALPH Loop}
\date{}
\maketitle

\begin{abstract}
This paper presents D7, formalizing the underlying principles using information-theoretic methods.

\textbf{Keywords}: D7, information theory, formalization.
\end{abstract}

\tableofcontents
\section{D7}




\subsection{0.}


\textbf{**Bayesian}**frequentist

\textbf{D7 }\textbf{coarse-graininginformation erasure}



\subsection{1.}


\subsubsection{1.1}


\textbf{ D7-1} \(\phi_t: M \to M\)  \(M\)

\textbf{ D7-2} \(\Gamma: M \to \mathcal{P}\) \(\mathcal{P}\)

****
\[\Delta S_{\text{erasure}}(t) = S[\Gamma(\mu_t)] - S[\Gamma(\mu_0)] \geq 0\]

\(\mu_t = \phi_{t*}(\mu_0)\)  Liouville


\subsubsection{1.2}


\textbf{ D7-1} \(A \subset M\) ****
\[\mathbb{P}_{\Gamma}[A] = \lim_{T \to \infty} \frac{1}{T} \int_0^T \mathbf{1}_A(\phi_t(x)) \, d\mu_{\Gamma}(x)\]

\(\mu_{\Gamma}\)

\textbf{*}* \(\Gamma\)  \(\phi_t\) **


\subsection{2.}


\subsubsection{ D7-1}


****
\((M, \phi_t, \mu)\) \(\Gamma\)
\[\mathbb{P}_{\Gamma} \longleftrightarrow Q_{\Gamma} = \inf_{t \geq 0} \left| I(X_0 : X_t | \Gamma) \right|\]

\(I(X_0 : X_t | \Gamma)\)  \(\Gamma\) \(Q_{\Gamma}\) ****

****


\subsubsection{ D7-2}


****
\[\frac{d}{dt} \mathbb{H}[\mathbb{P}_t] = -\Lambda(\Gamma, \phi_t)\]


\begin{itemize}
\item \(\mathbb{H}[\mathbb{P}_t]\)  \(\mathbb{P}_t\)
\item \(\Lambda(\Gamma, \phi_t)\) \textbf{Kolmogorov-Sinai }
\end{itemize}


\textbf{ D7-2a}\(\Lambda > 0\)****ergodic invariant measure

\textbf{ D7-2b}\(\Lambda = 0\)  ****


\subsubsection{ D7-3 =}


****
\(\mathbb{P}(\theta|x) \propto \mathbb{P}(x|\theta) \cdot \mathbb{P}(\theta)\)
\[\delta \mu = -\Gamma^{-1} \cdot \nabla_{\Gamma} \log \mathbb{P}(x|\Gamma)\]

\textbf{**""}**


\subsubsection{ D7-4 =}


****
\(X_i = f \circ \phi^i(x)\)
\[\frac{1}{n} \sum_{i=1}^{n} X_i \xrightarrow{a.s.} \int_M f \, d\mu_{\text{inv}}\]

\(\inf_{t} \left| I(X_0 : X_t | \Gamma) \right| > 0\)

\textbf{**""}**


\subsubsection{ D7-5}


****
\(\rho(t) = e^{-iHt/\hbar} \rho(0) e^{iHt/\hbar}\)
\[p_i = \langle i | \rho | i \rangle = \text{Tr}(\Pi_i \rho)\]


\[\mathbb{P}_{\text{quantum}} = \mathbb{P}_{\text{dynamical}} + \epsilon_{\text{pure-erasure}}\]

\(\epsilon_{\text{pure-erasure}}\) **** \(\phi_t\)  \(\{|i\rangle\}\)

\textbf{**  }**


\subsection{3.}


\subsubsection{ D7-A}


\textbf{ D7-2  D7-4}
1.  Lorenz/Chirikov  \(\Lambda_{KS} > 0\)
2.  \(\epsilon\)
3.  \(\propto \Lambda_{KS}(\epsilon)\)

**** \(\Lambda_{KS}\)  D7-2


\subsubsection{ D7-B}


\textbf{ D7-2b}
1. Harmonic oscillator, Toda lattice
2.
3. \(I(X_0 : X_t | \Gamma)\)  \(> 0\)

****


\subsubsection{ D7-C}


\textbf{ D7-3}
1.  SDE \(\dot{x} = f(x) + \sigma \xi(t)\)
2.  \(I(X_0 : \text{observations} | \Gamma)\)
3.  \(\approx\)


\subsubsection{ D7-D-}


\textbf{ D7-5}
1.
2.
3.  \(\mathbb{P}_{\text{dynamical}} + \epsilon_{\text{pure}}\)


\subsection{4.}


\subsubsection{4.1 D9}


| D9  | D7  |
|---------|---------|
| C3 |  \(Q_{\Gamma}\) |
| C4 |  |
| C5 |  |

\textbf{D7-D9 }Agent  \(\Gamma^*\)  \(Q_{\Gamma^*}\)


\subsubsection{4.2 S18}


\textbf{S18-}
\[\mathcal{F}_{\text{prob}}[\mathbb{P}, \Gamma] = \mathbb{H}[\mathbb{P}] + \lambda \cdot Q_{\Gamma}[\phi]\]

**** \(\mathbb{P}^* = \arg\min_{\mathbb{P}} \mathcal{F}_{\text{prob}}[\mathbb{P}, \Gamma]\)


\subsubsection{4.3 Chimera}


\textbf{D7-Chimera }Chimera  MPI > 0.8 \textbf{*}***


\subsection{5.}


\subsubsection{5.1}


1. \textbf{D7-1 } \(Q_{\Gamma}\)
2. \textbf{D7-2 }
3. \textbf{D7-5 }+

\subsubsection{5.2}


1. ****
2. ****
3.  \(\Gamma\) ""


\subsection{6.}


\#\#\#

**** \((M, \phi_t, \mu)\)  \(\Gamma: M \to \mathcal{P}\)

****
\[\boxed{\mathbb{P}_{\Gamma}[A] \equiv \mathbb{E}_{\mu}\left[\mathbf{1}_A \mid \mathcal{F}_{\Gamma}\right] = \lim_{T \to \infty} \frac{1}{T} \int_0^T \mathbf{1}_A(\phi_t(x)) \, d\mu(x)}\]

****
\[\frac{d}{dt} \mathbb{H}[\mathbb{P}_t] = -\Lambda_{KS}(\Gamma, \phi_t) \leq 0\]

****


*D7 *

\begin{thebibliography}{99}
\bibitem{ref1} Shannon, C.E. (1948). A Mathematical Theory of Communication.
\bibitem{ref2} Cover, T.M. (2006). Elements of Information Theory.
\end{thebibliography}

\end{document}